\section{Data and Method}\label{sec:approach}

This is a measurement-based project. Below we describe the data source, the basket
and country list we measured, the pipeline we built, and how the raw data becomes
the results in Section~\ref{sec:results}.

\subsection{Data Source}

The core measurement surface is public and queryable from any vantage point, which
is a significant practical advantage: we did not need machines or volunteers inside
censored countries. The public Apple iTunes Lookup API accepts a \texttt{country}
parameter, so the same app identifier can be queried against every national
storefront. A lookup-by-ID returns \texttt{resultCount}${\geq}1$ where the app is
available in that storefront and \texttt{resultCount}${=}0$ where it is not. This
is a cleaner signal than keyword search, which can return near-matches or
copycats. For ground truth we used documented removal events (the 2024 Russian VPN
purge, the 2017 China VPN and \emph{New York Times} removals) and Freedom House's
\emph{Freedom on the Net} 2024 per-country scores~\cite{fotn2024}.

We restricted this study to the Apple App Store. Google Play exposes no equally
clean per-country availability API; robust Play measurement requires scraping that
we judged too unreliable to complete within our timeframe, so we leave it to future
work and treat our findings as Apple-specific (Section~\ref{sec:conclusion}).

\subsection{App Basket and Country List}

We assembled a basket of \textbf{32 apps} across five categories
(the full list appears in Figure~\ref{fig:heatmap}): nine \textbf{VPNs} (e.g., NordVPN,
Proton VPN, Psiphon, Onion Browser), six \textbf{secure messengers} (Signal,
Telegram, WhatsApp, Threema, Session, Wire), six \textbf{news} apps (NYT, BBC, The
Guardian, Al Jazeera, Deutsche Welle, Voice of America), five \textbf{topically
sensitive} apps (Grindr and HER for LGBTQ+, Quran Majeed, a Dalai Lama app, and
the Bridgefy offline mesh messenger), and six \textbf{neutral control} apps
(Candy Crush, Duolingo, Google Maps, Google Translate, The Weather Channel,
Microsoft Word) that no government is expected to target. Every app identifier was
validated against the US baseline before inclusion: an ID that did not resolve in
the US store was corrected or dropped.

We measured \textbf{14 country storefronts}: known censors China, Russia, Vietnam,
Saudi Arabia, Pakistan, Egypt, the UAE, Turkey, and India, and free-expression
baselines the United States, France, Germany, the United Kingdom, and Japan. All 14
have a \emph{Freedom on the Net} 2024 score, spanning 9 (China) to 78 (UK and
Japan). We originally included Iran, but every Iranian query returned HTTP 400:
Apple operates \emph{no} storefront in Iran under US sanctions, so there is nothing
to query. This is a distinct mechanism from per-app removal, and we exclude Iran
from the measurable set and report it as a limitation.

\subsection{Pipeline}

We built a small Python pipeline (\texttt{collect.py}) that queries the full
$32 \times 14 = 448$-cell cross product, classifies each cell as \emph{available},
\emph{unavailable}, or \emph{error}, and \textbf{caches the raw JSON response for
every query} so the entire analysis is reproducible offline. The collector throttles
requests, retries transient failures with backoff, and keeps \emph{error} distinct
from \emph{unavailable} so coverage is reported honestly. Analysis
(\texttt{analyze.py}) uses Pandas and Matplotlib. All code, the raw response cache,
and the availability matrix are released openly.
